\documentclass[10pt,a4paper]{amsart}
\usepackage[margin=19mm]{geometry}
\usepackage{amsmath,amssymb,mathtools}
\usepackage[scr=rsfs,cal=euler]{mathalfa}
\usepackage{xfrac}
\usepackage[unicode,hidelinks,bookmarks=false]{hyperref}
\newcommand{\ie}{i.e.\ }
\newcommand{\cf}{cf.\ }
\newcommand{\resp}{respectively\ }
\newcommand{\Subsection}{\subsection}
\providecommand{\nobreakdash}{\nobreak\hbox{-}}
\setlength{\emergencystretch}{2em}
\multlinegap0pt
\usepackage{bm}
\usepackage{xcolor}
\usepackage{amssymb}
\usepackage{tikz}
\newtheorem{prop}{Proposition}
\title{Tropical curves with parallel rays}
\author{Song JuAe}
\date{Source: arXiv:2605.16910v1 (CC BY 4.0)}
\begin{document}
\maketitle

\noindent\emph{Source and licence.} Tropical curves with parallel rays. Version arXiv:2605.16910v1 (CC BY 4.0), distributed by arXiv under the Creative Commons Attribution 4.0 International licence, http://creativecommons.org/licenses/by/4.0/, as recorded in the arXiv licence metadata for this exact version and shown on the abstract page https://arxiv.org/abs/2605.16910v1. Full licence notice: \texttt{licenses/arxiv-2605.16910-CC-BY-4.0.txt}.\par\smallskip

\noindent\textit{Editorial scope.} Counted unit: the article's prop prop4-6-1 (source0.tex lines 1336--1346). The article's complete original proof of it is delivered with it (source0.tex lines 1348--1377, one \texttt{proof} environment of 30 lines). No mathematics is written, repaired or re-derived: the statement and the proof are the article's own lines, in the article's own notation, and no retained line is substituted or re-flowed.  No further source-local context is needed: every cross-reference of every retained line resolves inside the two delivered spans.  Nothing was omitted from any retained span, so the package carries no omission record.  No retained line contains a citation, so the wrapper carries no bibliography and no bibliography entry.  No retained line contains a figure environment, an image-inclusion command or drawing code, so no image is delivered and the package declares no asset dependency.  Editorial formatting only, in addition: an AMS \texttt{amsart} preamble that restates the article's own declarations and macros used by the retained lines, namely the environments \texttt{prop} on separate counters, together with the standard packages those lines need.  The retained environments print in this document as Proposition 1 in order of appearance, whereas the article prints the same items with the numbering of its own full document.  The complete native source of the article as distributed in the arXiv e-print for this exact version is bundled unchanged as \texttt{sources/200754/source0.tex}.
\begin{prop}
    \label{prop4-6-1}
The above $M_n$ is a semifield over $\boldsymbol{T}$.
Moreover, the map
\begin{align*}
\psi :\, \boldsymbol{T}[\boldsymbol{X}^{\pm}]_n \to M_n;
\bigoplus_{\boldsymbol{i} \in \boldsymbol{Z}^n} a_{\boldsymbol{i}} \odot \boldsymbol{X}^{\odot \boldsymbol{i}} \mapsto \boxplus_{\boldsymbol{i} \in \boldsymbol{Z}^n} a_{\boldsymbol{i}} \odot \boldsymbol{X}^{\odot \boldsymbol{i}},
-\infty \mapsto -\infty
\end{align*}
is a surjective $\boldsymbol{T}$-algebra homomorphism.
\end{prop}

\begin{proof}
The first assertion is straightforward.
The identity with respect to $\boxplus$ is the real number zero, and that with respect to $\boxdot$ is $-\infty$.
For $a \odot \boldsymbol{X}^{\odot \boldsymbol{i}} \in M_n \setminus \{ -\infty \}$, its multiplicative inverse is $a^{\odot (-1)} \odot \boldsymbol{X}^{\odot (-\boldsymbol{i})}$.

By definition, $\psi(0) = 0$ and $\psi(-\infty) = -\infty$ hold.
Let $f = \bigoplus_{\boldsymbol{i} \in \boldsymbol{Z}^n} a_{\boldsymbol{i}} \odot \boldsymbol{X}^{\odot \boldsymbol{i}}, g = \bigoplus_{\boldsymbol{i} \in \boldsymbol{Z}^n} b_{\boldsymbol{i}} \odot \boldsymbol{X}^{\odot \boldsymbol{i}} \in \boldsymbol{T}[\boldsymbol{X}^{\pm}]_n \setminus \{ -\infty \}$.
Here, all $a_{\boldsymbol{i}}$ and $b_{\boldsymbol{j}}$ except for a finite number of $\boldsymbol{i}, \boldsymbol{j} \in \boldsymbol{Z}^n$ are $-\infty$.
Since $a_{\boldsymbol{i}} \odot \boldsymbol{X}^{\odot \boldsymbol{i}} \oplus b_{\boldsymbol{i}} \odot \boldsymbol{X}^{\odot \boldsymbol{i}} = a_{\boldsymbol{i}} \odot \boldsymbol{X}^{\odot \boldsymbol{i}} \boxplus b_{\boldsymbol{i}} \odot \boldsymbol{X}^{\odot \boldsymbol{i}}$ holds by definition, 
\begin{align*}
\psi(f \oplus g) &= \psi \left( \bigoplus_{\boldsymbol{i} \in \boldsymbol{Z}^n} (a_{\boldsymbol{i}} \oplus b_{\boldsymbol{i}}) \odot \boldsymbol{X}^{\odot \boldsymbol{i}} \right)\\
&
= \boxplus_{\boldsymbol{i} \in \boldsymbol{Z}^n} (a_{\boldsymbol{i}} \oplus b_{\boldsymbol{i}}) \odot \boldsymbol{X}^{\odot \boldsymbol{i}}\\
&= \boxplus_{\boldsymbol{i} \in \boldsymbol{Z}^n} (a_{\boldsymbol{i}} \odot \boldsymbol{X}^{\odot \boldsymbol{i}} \oplus b_{\boldsymbol{i}} \odot \boldsymbol{X}^{\odot \boldsymbol{i}})\\
&= \boxplus_{\boldsymbol{i} \in \boldsymbol{Z}^n} (a_{\boldsymbol{i}} \odot \boldsymbol{X}^{\odot \boldsymbol{i}} \boxplus b_{\boldsymbol{i}} \odot \boldsymbol{X}^{\odot \boldsymbol{i}})\\
&= \left( \boxplus_{\boldsymbol{i} \in \boldsymbol{Z}^n} a_{\boldsymbol{i}} \odot \boldsymbol{X}^{\odot \boldsymbol{i}} \right) \boxplus \left( \boxplus_{\boldsymbol{i} \in \boldsymbol{Z}^n} b_{\boldsymbol{i}} \odot \boldsymbol{X}^{\odot \boldsymbol{i}} \right)\\
&= \psi(f) \boxplus \psi(g)
\end{align*}
hold.
Also for $c_{\boldsymbol{k}} := \bigoplus_{\boldsymbol{i}, \boldsymbol{j} \in \boldsymbol{Z}^n :\, \boldsymbol{i} \odot \boldsymbol{j} = \boldsymbol{k}} a_{\boldsymbol{i}} \odot b_{\boldsymbol{j}}$,
\begin{align*}
\psi(f \odot g) &= \psi \left( \bigoplus_{\boldsymbol{k} \in \boldsymbol{Z}^n} c_{\boldsymbol{k}} \odot \boldsymbol{X}^{\odot \boldsymbol{k}} \right)\\
&= \boxplus_{\boldsymbol{k} \in \boldsymbol{Z}^n} c_{\boldsymbol{k}} \odot \boldsymbol{X}^{\odot \boldsymbol{k}}\\
\psi(f) \boxdot \psi(g) &= \left( \boxplus_{\boldsymbol{i} \in \boldsymbol{Z}^n} a_{\boldsymbol{i}} \odot \boldsymbol{X}^{\odot \boldsymbol{i}} \right) \boxdot \left( \boxplus_{\boldsymbol{j} \in \boldsymbol{Z}^n} b_{\boldsymbol{j}} \odot \boldsymbol{X}^{\odot \boldsymbol{j}} \right)\\
&= \boxplus_{\boldsymbol{k} \in \boldsymbol{Z}^n} c_{\boldsymbol{k}} \odot \boldsymbol{X}^{\odot \boldsymbol{k}},
\end{align*}
and hence $\psi(f \odot g) = \psi(f) \boxdot \psi(g)$ holds.
Clearly $m \in M_n \subset \boldsymbol{T}[\boldsymbol{X}^{\pm}]_n$ goes to $m$ by $\psi$.
In conclusion, $\psi$ is a surjective $\boldsymbol{T}$-algebra homomorphism.
\end{proof}

\end{document}
